\section{Scope, contributions, and non-claims}
\label{sec:scope}

This work is best read as a \emph{proposal of a substrate model plus a diagnostic program}.
It is not yet a closed derivation of all Standard Model parameters.

\subsection{Contributions}
\begin{enumerate}
\item \textbf{Continuum substrate formulation.}
A rotationally invariant continuum model is stated as a hyperelastic action on a 3D reference domain,
with an embedding into $\R^4$ (\Cref{sec:continuum-model}).

\item \textbf{Geometric-phase diagnostics for gauge structure.}
A clear emergence map is given from a narrowband, polarized substrate wave to a
Berry/Pancharatnam-Berry connection and curvature, including the non-Abelian (Wilczek--Zee)
case for near-degenerate mode subspaces (\Cref{sec:geophase}).

\item \textbf{Localized modes as operator problems.}
A framework is provided to pose ``particle-like'' localized modes as self-adjoint eigenproblems
around spherically symmetric background configurations, avoiding ad-hoc periodicity or boundary
assumptions (\Cref{sec:localized}).

\item \textbf{Validation roadmap.}
A set of concrete, simulation-independent checks is listed (symmetry tests, dispersion,
holonomy extraction, gauge-invariant loop integrals), with discrete implementation moved to an
appendix (\Cref{sec:validation}, \Cref{app:discrete}).
\end{enumerate}

\subsection{Non-claims (what is \emph{not} asserted here)}
\begin{enumerate}
\item We do not claim a complete derivation of the Standard Model, nor a replacement for QED/QFT.
\item We do not claim that the Berry connection \emph{is} electromagnetism; rather, we propose it as a
candidate effective gauge variable whose field-like behavior must be empirically or numerically tested.
\item We do not claim a final, unique hyperelastic constitutive law; the paper isolates minimal
requirements (isotropy, stability, existence of a linear-wave regime) and keeps the constitutive choice modular.
\item We do not claim that the discrete lattice approximation is fundamental; it is treated as a numerical
realization of the continuum model and kept out of the main narrative.
\end{enumerate}
